\chapter{Introduction}\label{ch:intro}

\section{Motivation}
\label{sec:motivation}

A network dataplane is the part of a protocol stack that moves bytes on the hot path: parsing, validating, filtering, routing, buffering, and exchanging data with the transport. In the construction of network dataplanes, two notions organize the work. An \emph{atom} is the smallest useful component: a pure or effectful transformation with no hidden state, total on its declared domain. A \emph{molecule} is a composition of atoms (sequential or tensorial) with bounded internal state; a ring buffer, a parser, and the whole receive path are molecules. The engineering target is concrete: \emph{zero allocation on the hot path}: no heap traffic, no per-packet allocation, no copy that can be fused away.

The target is hard because the requirement is global (an allocation anywhere in a composed pipeline defeats the guarantee everywhere) and because the usual answers to compositionality, abstraction layers, dynamic dispatch, managed containers, are themselves the most common sources of allocation. A dataplane that must be allocation-free by construction must therefore be \emph{optimized by composition}: the object of study is a family of stateful transformations together with the rules for joining them, and the performance of a pipeline is a property of how its parts are wired.

Sequential machines are the natural model: Mealy machines \citep{mealy} compute deterministically from a fixed state space, and Moore's analysis \citep{moore} established the basic shape of minimization. What is missing is a compositional, algebraic treatment: a category whose objects are types and whose morphisms are machines, closed under the two composition schemes a dataplane uses (sequential ($\circ$) and parallel ($\otimes$)) and in which engineering questions have theorem-shaped answers. Does normalization preserve behavior and cost? Is the minimal state space computable? When does a loop provably converge? When is a pipeline provably allocation-free? This thesis answers them by developing the category \textbf{Mol} of stateful transformations and proving the structure theorems that make it a foundation.

\section{Thesis statement}
\label{sec:thesis-statement}

\textbf{Thesis.} The category \textbf{Mol} of stateful transformations (deterministic Mealy machines $(S, \mathit{step})$ with $\mathit{step} : S \times A \to B \times S$, quotiented by the state-space bijection condition) is a complete algebraic foundation for the construction of zero-allocation network dataplanes: every dataplane artifact is a molecule, molecules compose sequentially and in parallel within a symmetric monoidal category, and the engineering properties of dataplanes (normalization, minimization, loop convergence, batching, zero allocation, solvability of design problems) are theorems of Mol.

The claim rests on three supports. \emph{Compositionality:} Mol is symmetric monoidal \citep{maclane}, with tensor the tuple product, unit $()$, and coherence from Mac Lane's theorem; results about pipelines are results about morphisms. \emph{Completeness:} the catalog \autoref{thm:1}--\autoref{thm:68} covers the category structure (\autoref{thm:1}--\autoref{thm:8}, \autoref{thm:67}, \autoref{thm:68}), equational theories and rewrite systems (\autoref{thm:9}--\autoref{thm:19}, \autoref{thm:54}), free molecules and normal forms (\autoref{thm:20}--\autoref{thm:24}), cost enrichment and optimality (\autoref{thm:25}--\autoref{thm:32}, \autoref{thm:53}), pure/effectful decomposition (\autoref{thm:33}--\autoref{thm:37}), minimal state spaces (\autoref{thm:38}--\autoref{thm:41}), loops and feedback (\autoref{thm:42}--\autoref{thm:45}), batching (\autoref{thm:46}--\autoref{thm:48}), zero allocation (\autoref{thm:49}--\autoref{thm:50}, \autoref{thm:55}), design decision problems (\autoref{thm:51}--\autoref{thm:52}), and kernel-boundary optimization (\autoref{thm:56}--\autoref{thm:66}). \emph{Honesty:} every theorem carries a complete proof; [N] statements are original and proved with elementary means, [C] statements restate classical results with provenance cited; computational claims are verified by executable Rust tools whose logs appear in \autoref{app:a}; no conjecture is presented as a theorem.

\section{Contributions and theorem catalog}
\label{sec:contributions}

The theorems of this thesis are grouped by the chapter that states and
proves each. Statements without a citation are original to this thesis;
classical results are restated and proved inside Mol with the source
cited. Around the theorems stand the framework (\autoref{ch:mol}),
which recasts Mealy-style sequential machines \citep{mealy,moore} as
morphisms, the verification of computational claims by executable Rust
tools (logs in \autoref{app:a}), and a companion software standard that
turns the theorems into enforceable engineering policy.

\autoref{ch:mol} proves the category structure: composition,
well-definedness of the state-space quotient, symmetric monoidal
structure, the pure and effectful subcategories, bisimulation
congruence, tensor closure, and determinism. \autoref{ch:lawvere}
models rings, parsers, and transports as equational theories.
\autoref{ch:rewrite} proves termination, confluence, and Knuth--Bendix
completion for the rewrite engine. \autoref{ch:free} constructs free
molecules and proves finiteness of normal forms and their completeness
as a lookup table. \autoref{ch:enrich} proves cost optimality by
shortest paths and the refinement order. \autoref{ch:kleisli} proves
the decomposition of every hybrid molecule into pure, effectful, and
pure factors, and the traced structure of loops. \autoref{ch:coalgebra}
proves minimal state spaces and compositional bounds. \autoref{ch:trace}
proves contraction bounds for fixed-iteration loops. \autoref{ch:operad}
proves the batching laws. \autoref{ch:linear} proves cut elimination as
pipeline fusion and linearity as a zero-allocation guarantee.
\autoref{ch:situations} proves the finiteness and decidability of
design decision problems. \autoref{ch:breakthroughs} proves the
kernel-boundary laws, the affine crossover, the optimization lattice,
the band contraction of the admission estimate, and the route
automaton theorems.
\section{Scope}
\label{sec:scope}

This thesis is a synthesis-and-application thesis with deliberately bounded scope; the boundaries are part of the argument.

\textbf{In scope.} Mol and its structure theorems \autoref{thm:1}--\autoref{thm:68}, each with a complete proof; the application of the theory to dataplane engineering: normalization, dispatch, batching, minimization, loop analysis, zero-copy reasoning, kernel-boundary cost, and design decision problems; and the verification of every computational claim by executable Rust tools, with PASS logs in \autoref{app:a} and catalogs in \autoref{app:b}.

\textbf{Out of scope.} This thesis does not construct a transport engine. A companion implementation of TCP, UDP, and SCTP lives in the same repository, in the \texttt{Code/} tree, and is not a substitute for the proofs. Machine-specific performance techniques (socket options, polling strategies, SIMD, syscall catalogs) appear only in the companion standard's non-normative appendix, not as theorems. A complete machine-checked formalization of the catalog is future work and is not claimed.

\textbf{Methodological filter.} No conjecture is presented as a theorem; any item that cannot be proved as stated is weakened, stated precisely, or moved to future work. [C] statements restate classical results proved inside Mol, with the original cited; [N] statements are original and proved with elementary machinery. Every computational claim is discharged by an executable Rust tool; no argument relies on unverified computation.

\section{Roadmap}
\label{sec:roadmap}

\textbf{Part I: Foundations} (\autoref{ch:prelim}--\autoref{ch:kleisli}). \autoref{ch:prelim} fixes the categorical machinery: monoidal, symmetric monoidal, and traced monoidal categories \citep{maclane,joyalsv}; monads and Kleisli; coalgebras; Lawvere theories \citep{lawvere}; operads and clubs \citep{leinster,kelly}; effectus theory; linear logic \citep{girard}; locally presentable categories \citep{adamek}. \autoref{ch:mol} defines Mol and its subcategories (\autoref{thm:1}--\autoref{thm:8}); \autoref{ch:lawvere} models rings, parsers, and transports as Lawvere theories (\autoref{thm:9}--\autoref{thm:12}); \autoref{ch:rewrite} develops the rewrite engine, with Knuth--Bendix completion \citep{knuthbendix} and Newman confluence \citep{newman} (\autoref{thm:13}--\autoref{thm:19}, \autoref{thm:54}); \autoref{ch:free} constructs free molecules and normal forms (\autoref{thm:20}--\autoref{thm:24}); \autoref{ch:enrich} proves cost optimality and refinement (\autoref{thm:25}--\autoref{thm:32}, \autoref{thm:53}); \autoref{ch:kleisli} proves the decomposition $M = q\circ e\circ p$ (\autoref{thm:33}--\autoref{thm:37}).

\textbf{Part II: Structure theorems} (\autoref{ch:coalgebra}--\autoref{ch:linear}). \autoref{ch:coalgebra}: minimal realization and compositional bounds \citep{nerode,hopcroft} (\autoref{thm:38}--\autoref{thm:41}). \autoref{ch:trace}: feedback and fixed-iteration loops, with Banach--Dobrushin contraction bounds \citep{banach,dobroushin} (\autoref{thm:42}--\autoref{thm:45}). \autoref{ch:operad}: operadic composition and batching (\autoref{thm:46}--\autoref{thm:48}). \autoref{ch:linear}: cut elimination as fusion, linearity as zero allocation (\autoref{thm:49}--\autoref{thm:50}, \autoref{thm:55}).

\textbf{Part III: Applications} (\autoref{ch:situations}--\autoref{ch:conclusion}). \autoref{ch:situations}: a situation is a triple $(A\to B, \Phi, \bar{c})$, a solution a molecule $M$ with $M \vdash \Phi$ and $\mathrm{cost} (M) \leq \bar{c}$ (\autoref{thm:51}--\autoref{thm:52}). \autoref{ch:testing}: law tests, bisimulation \citep{park}, model-based testing, fuzzing, cost regression. \autoref{ch:related}: related work. \autoref{ch:conclusion}: summary, limitations, future work. \autoref{app:a}: proof-checking scripts and PASS logs; \autoref{app:b}: optimization and testing catalogs.
